\originalchapter{均匀膨胀、能量积分与尘埃宇宙习题}
本章包括8.286 PS3四题与PS4五题。Newton模型的有限球或无限圆柱是受控数学模型；GR的Friedmann方程另有压强项。以下每次调用守恒能量均说明它属于哪个系统，避免把局部第一定律、类比粒子能量和整宇宙能量混为一谈。

\src{8.286 PS3.1，圆柱宇宙，(a)--(e)。}
\begin{exercise}
无限长均匀圆柱初半径$R_{\max,i}$、密度$\rho_i$，只在横向膨胀，沿$z$无运动。初$t_i$各壳速度$\dot r_i=H_i r_i$。
(a)由引力Gauss律求$g=-A\mu(r)/r$的常数，$\mu$是包围线质量；
(b)给初半径$r_i$的壳写$\ddot r$；
(c)用$u=r/r_i$证明均匀膨胀、定义$a(t)$；
(d)求$\rho(a)$、$\ddot a$；
(e)求$\tfrac12\dot a^2+V(a)$守恒量，判是否永久膨胀。
\end{exercise}
\begin{solution}
长度$L$的共轴Gauss柱面有$2\pi rL\,g=-4\pi G\mu L$，故$A=2G$。无壳交叉时每壳内部线质量不变，$\mu=\pi r_i^2\rho_i$，所以
\[
 \ddot r=-\frac{2\pi G\rho_i r_i^2}{r},\qquad
 \ddot u=-\frac{2\pi G\rho_i}{u},\quad
 u(t_i)=1,\ \dot u(t_i)=H_i.
\]
方程与初条件均不依赖$r_i$，在$u>0$处唯一解相同，故$u=a$。横向面积增为$a^2$，长度不变，$\rho=\rho_i/a^2$，$\ddot a=-2\pi G\rho a$。乘$\dot a$积分得
\[
 \mathcal E=\frac12\dot a^2+2\pi G\rho_i\ln a,\qquad
 V(a)=2\pi G\rho_i\ln a.
\]
对任何有限初膨胀速度，$V\to+\infty$禁止$a\to\infty$，最大值$a_{\max}=\exp[H_i^2/(4\pi G\rho_i)]$后收缩。与三维球的$-1/a$势不同，圆柱的对数势不能允许有限能量逃逸。这是横向Newton模型，不是各向同性FLRW宇宙。
\end{solution}

\src{8.286 PS3.2，$a=bt^{3/4}$，(a)--(b)。}
\begin{exercise}
平坦宇宙的特殊物质使$a=bt^{3/4}$，求(a)$H(t)$、(b)$\rho(t)$。可用所有物质均适用的$H^2=8\pi G\rho/3$，不能用仅尘埃适用的$\ddot a=-(4\pi G/3)\rho a$。
\end{exercise}
\begin{solution}
直接对数微分$H=3/(4t)$，从约束方程得
\[
 \rho=\frac{3H^2}{8\pi G}=\frac{27}{128\pi Gt^2}.
\]
常数$b$消失。此$a$不对应无压尘埃；若进一步使用能量守恒，$\rho\propto a^{-8/3}$意味着$p=-\rho c^2/9$，负压正是尘埃加速度方程不适用的原因。
\end{solution}

\src{8.286 PS3.3，Friedmann积分常数的能量意义，(a)--(c)。}
\begin{exercise}
有限均匀尘埃球初半径$R_{\max,i}$、密度$\rho_i$，$a(t_i)=1$，有$\mathcal E=\tfrac12\dot a^2-(4\pi G\rho_i)/(3a)$。
(a)对初$r_i$、质量$m$的共动测试粒子，证明$m r_i^2\mathcal E$等于$\tfrac12mv^2-GmM(r_i)/r$，并说明这是沿壳运动的有效能量；
(b)逐壳积分总动能，求$K=c_K MR_{\max,i}^2\dot a^2/2$的$c_K$；
(c)逐壳装配求势能，证明相同系数给$E_{\rm total}=c_K MR_{\max,i}^2\mathcal E$。
\end{exercise}
\begin{solution}
(a)$r=ar_i,v=\dot a r_i,M(r_i)=4\pi\rho_i r_i^3/3$，直接相乘即给所求。球外壳在粒子位置处虽不产生力，却给随时间改变的势能常数；真势$\Phi(r,t)$对固定$r$的表达与这个只沿所选壳正确的$\Phi_{\rm eff}=-GM(r_i)/r$不同。因此“有效能量守恒”不等于任一粒子在时变全球势中的真实机械能守恒。

(b) 质量壳$\dd M=4\pi\rho_i r_i^2\dd r_i$，速度$\dot a r_i$，所以
\[
 K=\frac12\dot a^2\int_0^{R_{\max,i}}r_i^2\dd M
 =\frac12\dot a^2\frac35MR_{\max,i}^2,\quad c_K=\frac35.
\]
(c)在实际半径$R=aR_{\max,i}$逐壳带入质量，$\dd U=-GM(r)\dd M/r$，积分$U=-3GM^2/(5R)$。用$M=4\pi\rho_iR_{\max,i}^3/3$改写，
\[
 U=\frac35MR_{\max,i}^2\left(-\frac{4\pi G\rho_i}{3a}\right),
\]
故$c_U=c_K=3/5$。这里总能量属于无外力、无压、有限Newton球，不能拿来宣称一般无限FLRW时空具有同一个全局能量。
\end{solution}

\src{8.286 PS3.4，新增排斥力，(a)--(d)。}
\begin{exercise}
上述均匀球模型改为$\ddot r=-GM(r_i)/r^2+\gamma r^n$，$\gamma>0$不依赖质量。
(a)给$u=r/r_i$的方程；
(b)哪种$n$使它不依赖$r_i$；
(c)列保证统一解的初条件；
(d)对该$n$积分一次得到只含$a,\dot a$的能量关系。原题指出2011测验有官方解答；此处独立推导。
\end{exercise}
\begin{solution}
令$C=4\pi G\rho_i/3$，有$\ddot u=-C/u^2+\gamma r_i^{n-1}u^n$。非零$\gamma$时需$n=1$才对所有壳相同；初条件仍$u(t_i)=1,\dot u(t_i)=H_i$。故
\[
 \ddot a=-\frac C{a^2}+\gamma a,\qquad
 \mathcal E=\frac12\dot a^2-\frac C a-\frac{\gamma a^2}{2}=\text{常数}.
\]
排斥项对应负二次势，符号可由$-\dd V/\dd a$核查。它与宇宙常数的Newton极限形式相同，但给定的力模型本身还没有证明完整GR几何及能量守恒关系。
\end{solution}

\src{8.286 PS4.1，闭尘埃宇宙，(a)--(d)。}
\begin{exercise}
给$ct=\alpha(\eta-\sin\eta)$、$a/\sqrt{k}=\alpha(1-\cos\eta)$，$\alpha$为长度常数、$k>0$。
(a)求$H$；(b)求$\rho$；(c)由前两者求$\Omega=\rho/\rho_c$；
(d)求物理粒子视界，证明$\ell_{p,h}=\alpha\eta(1-\cos\eta)$。
\end{exercise}
\begin{solution}
令曲率半径$R=a/\sqrt{k}$，$\dd\eta/\dd t=c/[\alpha(1-\cos\eta)]$。因此
\[
 H=\frac{c\sin\eta}{\alpha(1-\cos\eta)^2},\qquad
 \rho=\frac{3c^2}{4\pi G\alpha^2(1-\cos\eta)^3}.
\]
后式由$H^2+c^2/R^2=8\pi G\rho/3$及$\sin^2\eta+(1-\cos\eta)^2=2(1-\cos\eta)$得到。再除$3H^2/(8\pi G)$，
\[
 \Omega=\frac{2}{1+\cos\eta}.
\]
最大膨胀处$H=0$，定义中的临界密度为零，$\Omega$发散，不表示$\rho$发散。径向光线的共形路程为$\int c\dd t/R=\eta$，所以按路径长定义$\ell_{p,h}=R\eta$。闭空间中若路程超过半周，最短两点距离与已走光路不再相同，视界已可能绕过空间；公式仍给光路积累长度。
\end{solution}

\src{8.286 PS4.2，开尘埃宇宙，(a)--(f)。}
\begin{exercise}
给$ct=\alpha(\sinh\eta-\eta)$、$a/\sqrt{\kappa}=\alpha(\cosh\eta-1)$，$\kappa=-k>0$。
(a)--(d)求$H,\rho,\Omega,\ell_{p,h}$；
(e)用早时级数求$a(t)$并估近似失效尺度$t_*$；
(f)求早时$1-\Omega$的幂律及系数。
\end{exercise}
\begin{solution}
与闭情形不同，约束为$H^2-c^2/R^2=8\pi G\rho/3$，结果
\[
 H=\frac{c\sinh\eta}{\alpha(\cosh\eta-1)^2},\quad
 \rho=\frac{3c^2}{4\pi G\alpha^2(\cosh\eta-1)^3},\quad
 \Omega=\frac{2}{1+\cosh\eta},\quad
 \ell_{p,h}=\alpha\eta(\cosh\eta-1).
\]
早时$\sinh\eta-\eta=\eta^3/6+O(\eta^5)$、$\cosh\eta-1=\eta^2/2+O(\eta^4)$，于是
\[
 \eta\simeq(6ct/\alpha)^{1/3},\qquad
 a(t)\simeq\frac{\sqrt{\kappa}\alpha}{2}(6ct/\alpha)^{2/3},
\]
\[
 1-\Omega=\tanh^2(\eta/2)\simeq\frac14(6ct/\alpha)^{2/3}.
\]
展开要求$\eta\ll1$，可取$t_*\sim\alpha/(6c)$为量级标尺；精确$\eta=1$给$t=\alpha(\sinh1-1)/c=0.1752\alpha/c$，并无唯一锐利“失效时刻”。早时$\Omega\to1$不代表空间拓扑或曲率符号改变，而是物质密度比曲率项增长得更快。
\end{solution}

\src{8.286 PS4.3，收缩期观察者，完整三个要求。}
\begin{exercise}
闭尘埃宇宙在收缩期$\eta_0>\pi$，观察者量得$H_0<0,\Omega_0>1$。用二者求距离大坍缩的剩余时间、视界内最负的蓝移$z$、以及该蓝移物体的回看时间。
\end{exercise}
\begin{solution}
由上一题对应闭解，先选\emph{收缩支}
\[
 \eta_0=2\pi-\arccos(2/\Omega_0-1),\qquad
 \frac{\alpha}{c}=\frac{\Omega_0}{2|H_0|(\Omega_0-1)^{3/2}}.
\]
若只用反三角函数的主值，会错误选择膨胀支。所求剩余时间为
\[
 \Delta t_{\rm crunch}=\frac{\alpha}{c}(2\pi-\eta_0+\sin\eta_0).
\]
任一发射时刻$1+z=a_0/a_e$，最蓝需要最大的$a_e$，即最大膨胀$\eta_e=\pi$。该事件的共形距离$\eta_0-\pi<\pi$确在 过去光锥 能达到的闭空间范围内。因此
\[
 z_{\min}=\frac{1-\cos\eta_0}{2}-1=-\frac1{\Omega_0},\qquad
 t_{\rm lookback}=\frac{\alpha}{c}(\eta_0-\sin\eta_0-\pi).
\]
更早发射处尺度因子更小，反而蓝移不再那么大，甚至可红移；“宇宙现在收缩”不意味着收到的所有远处光都蓝移。
\end{solution}

\src{8.286 PS4.4，$\Omega\to1^-$的年龄连续性。}
\begin{exercise}
开尘埃解的年龄表达式
\[
 Ht=\frac{\Omega}{2(1-\Omega)^{3/2}}
 \left[\frac{2\sqrt{1-\Omega}}{\Omega}
 -\operatorname{arsinh}\frac{2\sqrt{1-\Omega}}{\Omega}\right]
\]
在$\Omega\to1^-$时证明趋$2/3$。原题结尾另一句把$\Omega\to1^+$误印成$\Omega\to0^+$；此处只完成明确要求的左极限。
\end{exercise}
\begin{solution}
设$\delta=1-\Omega$，利用$\operatorname{arsinh}(2x/(1-x^2))=2\operatorname{artanh}x$，$x=\sqrt\delta$，可将式子精确改写为
\[
 Ht=\frac1\delta-\frac{1-\delta}{\delta^{3/2}}\operatorname{artanh}\sqrt\delta.
\]
展开$\operatorname{artanh}\sqrt\delta=\sqrt\delta+\delta^{3/2}/3+\delta^{5/2}/5+\cdots$，两个$1/\delta$项取消，余下
\[
 Ht=\frac23+\frac{2}{15}\delta+O(\delta^2)\longrightarrow\frac23.
\]
写出取消项是必要步骤；直接对极小$\delta$相减会损失数值精度。也可由开解$\eta\to0$直接取$\sinh\eta(\sinh\eta-\eta)/(\cosh\eta-1)^2$，同样得到$2/3$。
\end{solution}

\src{8.286 PS4.5，两个各向同性中心，加分题。}
\begin{exercise}
欧氏空间的物理分布关于两个不同点都各向同性，证明它均匀。原读物两圆相交论证没有覆盖不相交的情况；给出适用于全部距离的证明，可用已知欧氏几何事实。此断言不可直接推广到任意非欧氏空间。
\end{exercise}
\begin{solution}
设两中心$A,B$，距离$d>0$。在包含$AB$的二维平面内，取过两点、垂直该平面的平行轴作两个半周转动；它们的复合是沿$AB$的$2d$平移。二维时轴为平面法向，三维时也是真正的旋转，未偷用三维的点反演。每次转动都是给定的对称性，故这个平移也是。

再用绕$A$的任意旋转共轭，可得到任意方向、长度$2d$的平移。任意长度$\le4d$的向量都能写成两个长度$2d$向量的和，因此任意这样的短平移也保持分布。较长向量拆成有限个短向量，复合后得到任意平移对称。任意两点可被平移相连，物理性质处处相同，即均匀。

此证对维数至少2的欧氏空间成立，不需两圆相交。在球面上，关于南北两极旋转对称的纬度依赖分布可以仍不均匀，说明欧氏平移群是论证的实质。
\end{solution}
